\chapter{理论基础与文献综述}
%%====================================================================

\section{开源软件开发与开发者激励}

\subsection{开源软件的特征与发展历程}

开源软件（Open-Source Software, OSS）以源代码公开为基本特征，
允许任何人自由查看、使用、修改和分发\cite{lakhani2003open}。
有别于商业软件的封闭开发模式，开源项目在本质上依托互联网展开分布式协作——
参与者分布于全球各类机构，多数不受雇佣合同约束，
以自愿形式贡献代码、文档和讨论。协调工作则依靠版本控制系统、
邮件列表及代码审查流程等技术与社会机制得以维系。

自由软件运动的理念基础由 Richard Stallman\footnote{Richard Stallman 于1983年发起 GNU 项目，并于1985年创立自由软件基金会（Free Software Foundation, FSF），提出了"自由软件"的概念与 GPL 许可证体系。} 在20世纪80年代奠定；
1991年 Linus Torvalds 发布 Linux 内核，将该理念推向工程实践的前沿。
21世纪以来，互联网基础设施快速普及，加之 Git\footnote{Git 是 Linus Torvalds 于2005年为管理 Linux 内核源代码而开发的分布式版本控制系统，目前已成为软件工程领域的标准协作工具。} 版本控制工具日趋成熟，
全球开发者协作的组织成本大幅下降，以 GitHub 为核心的开源生态由此成形。
如今，Linux 内核支撑着全球绝大多数服务器与移动设备，
Python 已成为数据科学和人工智能研究的通用语言——
这些产品无一不是数以千计志愿开发者多年协作的成果。

从经济学角度审视，开源软件具备公共品的典型特征：
非竞争性与非排他性叠加，潜在的”搭便车”问题由此产生——
理论上人人都倾向于坐享他人贡献，而不愿自行承担成本\cite{lerner2002simple}。
但开源生态的实际状况与上述悲观预测相悖：
大量高质量项目维持了数十年的持续演进，贡献者遍布全球。
这种”反常”现象持续激发学界对开源开发者动机机制的探索\cite{gerosa2021shifting}，
并催生了一套有别于传统经济人假设的理论体系。

开源软件的经济重要性在近十余年间经历了质的跃升。
据估计，企业核心软件产品中超过70\%的代码来源于开源组件，
全球开源软件每年为商业应用贡献的价值已达数千亿美元规模。
然而，支撑这一巨大价值的开发者群体在规模上极度有限，
且整体处于志愿贡献状态——许多关键项目的核心维护工作仅集中在个位数的开发者身上\cite{lu2021core,jamieson2024predicting}。
“巨大经济价值与极少志愿维护者”之间的结构性不对称，
使理解和维持开发者参与动机成为兼具学术价值与现实紧迫性的议题。
近年来，”开源可持续性”这一概念被学界广泛用以描述上述挑战，
研究者呼吁在机构设计和激励机制层面作出系统性应对\cite{alami2024sustainability,wang2025opensource,qi2021opensource}。
本研究正是在该宏观背景下，对开源激励机制中尚未被充分探讨的
信息维度予以回应。

\subsection{开发者参与动机的理论框架}

开源开发者动机研究的核心挑战在于：为何大量技能成熟的开发者
愿意将大量时间投入缺乏直接经济回报的项目？
早期研究集中于内在动机的识别与刻画。
Shah（2006）通过对多个开源社区的质性考察指出，开发者参与的最初驱动力
往往源于对技术问题本身的强烈兴趣——即所谓”抓自己的痒”：
开发者为解决自身实际技术需求而参与，贡献本身即构成收益，
无需外部奖励加以维持\cite{shah2006motivation}。
Davidson 等人（2014）进一步确认，利他主义（即希望为社会留下有价值的成果）
在资深开发者中尤为突出，与兴趣、声誉追求一道构成多层次的内在动机结构\cite{davidson2014older}。

Lerner 和 Tirole（2002）从经济学视角提出了声誉机制假说：
开源社区中代码完全公开，每位开发者的贡献质量对所有人可见，
优秀贡献能够帮助开发者在同行社区中积累声誉，
而该声誉又可转化为劳动力市场上的职业优势\cite{lerner2002simple}。
按此逻辑，声誉实质上是一种”延迟货币”，
表面无偿的贡献可被解读为对人力资本的理性投资。
Hertel 等人（2003）基于对 Linux 内核贡献者的大规模调查，
将开发者动机归纳为实用主义需求、社区认同与政治信念三个维度，
着重论证了集体层面的心理认同对参与行为所具有的独立驱动作用\cite{hertel2003motivation}。
Fang 和 Neufeld（2009）的纵向研究则证实了职业发展动机的重要性——
将开源贡献视为职业发展手段的开发者，其持续参与的稳定性
显著高于纯粹基于兴趣的参与者\cite{fang2009understanding}。
晚近，Tulili 等人（2025）追踪开源生态系统中核心团队的演变过程，
系统刻画了核心开发者流入、流出与留存的动态模式，
确认维持核心团队稳定性的关键在于社区对成员贡献的可见化反馈机制\cite{tulili2025exploring}。

在内在动机与外在动机的理论谱系中，楼天阳等（2014）通过虚拟社区的实验研究报告了一项重要发现：
不同激励政策对成员参与动机的影响差异显著，
物质激励与社会激励的效果取决于社区成员的初始动机结构\cite{lou2014incentive}。
自我决定理论（Self-Determination Theory, Ryan \& Deci, 2000）\cite{ryan2000self}则提供了更为精细的分析框架，
区分了外在动机的不同内化程度：
一旦外在目标被个体充分内化——即个体认为追求该目标与其核心价值观一致——
外在动机便不再与内在动机对立，二者反而协同增强行为动机。
就开源社区而言，若开发者将接受外部资助理解为”对自身技术贡献价值的社会认可”
而非”为金钱而工作”，外在奖励并不必然触发挤出效应；
反之，若奖励被感知为外部控制手段，内在动机遭到侵蚀的可能性则大为增加。
这一分析暗示，捐赠信息的披露方式（究竟以”认可贡献价值”还是”购买服务”
的框架呈现）可能对最终激励效果产生关键影响\cite{benabou2003intrinsic}。

内在动机与外在动机并非简单相加，二者之间的交互关系颇为复杂，
学界以”挤出效应”概括之。
Bénabou 和 Tirole（2003）在理论模型中论证：当个体内在动机足够强时，
引入外在奖励反而可能适得其反——
外在奖励传递了”该工作本身不够有价值”的隐性信号，
同时破坏了个体基于内在动机行动所获得的社会认可\cite{benabou2003intrinsic}。
Qiao 等人（2021）在在线志愿知识贡献平台的实证研究中印证了这一机制：
货币激励短期内确能提升贡献量，
却同时显著降低了高内在动机用户的长期参与意愿\cite{qiao2021mitigating}。
王盼盼等（2022）针对在线健康社区的研究得到了类似结论——
有偿奖励短期内显著提升了医生的知识贡献数量，
但奖励机制撤除后，此前高度依赖内在动机参与的成员
贡献行为出现了显著回落\cite{wang2022reward}。
另一方面，社会性动机同样不可忽视。
Yang 等人（2021）区分了社会影响对初始参与与持续参与的差异化效应，
指出社会认同和群体规范主要驱动开发者作出初始贡献，
持续参与则更多依赖基于互动与信任建立起来的社区归属感\cite{yang2021differential}。
秦敏和李若男（2020）基于在线用户社区的研究同样确认，
社会影响和社区归属感是用户贡献行为形成的关键机制\cite{qin2020contribution}；
秦敏等（2015）更早期的研究则从复杂适应系统理论出发，
阐释了开放式创新社区中用户贡献行为的涌现机制\cite{qin2015cas}。
Sun 等人（2017）报告了货币奖励效果受社会连接性显著调节的证据：
社交关系紧密的平台上，货币奖励与社会认可的结合
能够产生超越单一维度激励的协同效果\cite{sun2017motivation}。
在加密货币生态这一特殊背景下，Kobayakawa 等人（2020）还注意到
代币市场价格表现通过影响开发者对项目前景的预期，
间接左右了开发者的参与强度，表明经济信号在特定情境下
可超越直接货币转移而发挥激励作用\cite{kobayakawa2020impact}。

\subsection{货币激励与捐赠赞助机制的研究综述}

尽管内在动机在开源参与中占据主导地位，随着开源软件战略重要性持续上升，
外部货币激励机制也逐步进入开源生态，引发学界的系统关注。
GitHub 于2019年推出 Sponsors 平台，允许个人和企业直接向开源开发者
支付定期赞助费用，由此构建了迄今最大规模的平台级经济支持基础设施。
Shimada 等人（2022）对该机制开展了最早的系统性研究，报告超过四分之三的
开发者在获得赞助后贡献量有所提升，且部分开发者将接受赞助本身
理解为对开源工作价值的认可，而非单纯的经济交换\cite{shimada2022github}。
Wang 等人（2022b）借助自然实验分析证实，赞助不仅正向影响受赞助开发者的贡献，
还通过社区观察机制对周围开发者产生有限的溢出激励效应\cite{wang2022influence}。
Zhang 等人（2022）考察了”谁赞助谁”的决定因素，确认赞助机制
更倾向于奖励已在社区中建立声誉的活跃开发者，因而更多承担
贡献后的事后认可功能，而非对潜在贡献的事前激励\cite{zhang2022who}。

不过，另一批研究呈现了截然不同的图景。
Overney 等人（2020）对多平台开源捐赠机制的大规模分析指出，
捐赠对开发者活动水平的整体提升作用十分有限——
多数项目即便获得捐赠，开发者的活动模式也并未发生显著变化\cite{overney2020not}。
Conti 等人（2023）追踪 GitHub Sponsors 数年的纵向数据后报告，
实际收到赞助的开发者在长期内创新型贡献有所减少，
经济激励可能将开发者的注意力从开放性技术探索
转向维持赞助关系的稳定性工作\cite{conti2023incentivizing}。
Wang 等人（2022a）在自然实验框架下则观察到显著的
“知识溢出挤出”效应：货币化激励使贡献者减少了免费分享知识的行为\cite{wang2022monetary}。

GitHub Sponsors 研究与非盈利基金会资助研究之间存在一个不容忽视的结构性差异。
GitHub Sponsors 属于”开发者对开发者”或”粉丝对开发者”的点对点直接支持机制，
赞助者自主选择受赞助对象，透明度和即时性均较高；
以 Brink 为代表的专项基金会则充当”中间人”角色，
先集中募资，再经申请审核程序将资金分配给具体开发者。
结构上的差异导致信息传播路径根本不同：
GitHub Sponsors 模式下赞助关系的建立本身即为双方直接互动，
赞助信号直达受赞助者，对旁观者亦具有一定的公开可见性；
基金会模式下，捐款信息的公开传播与资金的最终分配则完全分离——
公告受众远大于最终受益者，该结构天然形成了独立考察信息效应的实验条件。
正是”信息先于货币传播、信息受众大于货币受益者”这一特征，
赋予了基金会模式下信息披露研究超越 GitHub Sponsors 研究的独特价值。
代君等（2019）对开源项目协同开发行为特征的研究亦指出，
开发者之间的协同互动模式对项目成功至关重要，
外部信号对协同行为的激活效应需在结构化的协作网络中加以理解\cite{dai2019collaborative}。

综合以上研究，可识别出一个共同的局限：
现有文献对货币激励的理解普遍聚焦于实际货币转移——即开发者确实收到赞助金额这一事实，
进而将激励效应等同于直接经济补偿的作用。
但捐赠公告所携带的信息价值，
即资金实际到达之前甚至之后，信息本身对未直接受益的开发者所可能产生的激励效应，
在上述研究中均未被单独识别和检验。
捐赠信息披露作为独立于货币转移的激励机制，其有效性究竟如何？
这构成了本研究有别于既有文献的核心问题。

\section{信息披露的理论与实证}

\subsection{自愿信息披露理论}

信息披露研究的系统性发展始于会计学领域，后被管理学和信息系统学科广泛借鉴。
自愿信息披露理论（Voluntary Disclosure Theory）最早由 Verrecchia（1983）
在会计学框架中予以系统阐述\cite{verrecchia1983discretionary}。该理论的核心命题是：
信息不对称环境中，拥有私有信息的组织会权衡披露成本与收益，
自主决定是否以及如何向外部受众传递信息。
一旦披露预期带来的收益（如降低信息使用者的不确定性、
提升对组织的信任、减少信息使用者要求的风险溢价）超过披露成本，
组织便倾向于主动披露\cite{verrecchia1983discretionary}。

对非营利组织而言，信息披露具有维护公信力的特殊意义。
捐赠者和社会公众高度关注非营利组织的资金来源与使用情况，
透明的财务信息披露是此类组织维持捐赠关系、建立社会信任的基础性手段\cite{liu2015nonprofit}。
对 Bitcoin Brink 这类依赖社区捐赠运营的开源资助机构来说，
信息透明度不仅是合规性要求，更是与开发者社区建立信任关系的核心机制。
Brink 在 X 平台公开每笔收到的捐款，在官网公布每位受资助开发者名单——
这种主动且详细的信息披露行为，既是对捐赠者问责的承诺，
也是向开发者社区传递”项目获得外部认可”这一积极信号的渠道。
从自愿披露理论的视角审视，Brink 的披露行为本身即为一种精心设计的信息策略，
其对开发者行为的影响效果值得系统评估。

\subsection{信号理论}

信号理论（Signaling Theory）源自 Spence（1973）
对劳动力市场信息问题所做的经典分析\cite{spence1973job}，
此后被组织行为学、战略管理和信息系统等领域广泛采用。
该理论的基本洞察在于：信息不对称情境下，
处于信息劣势的一方（信号接收者）需依靠可观察的信号
推断自身无法直接观察的信息；
信息优势方（信号发送者）则可通过发送具有可信性的信号，
降低对方的不确定性并引导其作出有利的判断与行动。
有效信号须同时满足两个条件：
其一，与所传递的内容真实相关（相关性）；
其二，具有足够的成本使得虚假信号不值得发送（可信性）\cite{connelly2011signaling}。

在开源捐赠信息披露的情境中，Brink 的两种披露构成了性质迥异的信号。
非定向披露传递的是”项目获得社会认可与外部资金支持”的整体性信号，
类似于众筹平台上项目获得早期资助所产生的”热度信号”——
它向整个开发者群体提供了”该社区正在获得外界关注”的正面判断依据；
定向披露传递的则是”高质量贡献会获得认可与回报”的差异性信号，
通过明确揭示”谁因何贡献而获得资助”强化努力—奖励的因果链条，
向未入选者传递关于资助门槛和评判标准的具体信息，
进而影响其对”自身是否也有可能获得资助”的预期评估。
两类信号在内容、受众和激活机制上存在本质差异，
这构成了本研究分别检验其激励效应的逻辑起点。

\subsection{奖励信息披露与受众行为}

在信息披露对受众行为影响的实证研究中，众筹领域提供了与本研究最为接近的参照情境。
Mollick（2014）在众筹领域的开创性研究中确认，
项目信息的披露质量与透明度是决定众筹成败的关键因素\cite{mollick2014dynamics}。
Kim 等人（2022）进一步指出，风险信息的披露方式
显著影响投资者的参与决策，信息透明度的提升能有效降低感知风险\cite{kim2022risk}。
Conti 等人（2025）在创业融资情境中论证了另一条路径：
初创企业对开源社区的参与信息能向投资者传递创新能力信号，
信息透明度是连接资金提供方与接收方信任关系的核心机制\cite{conti2025beefing}。
Qiu 和 Zhang（2025）以精细加工可能性模型（ELM）为分析框架，
报告了众筹平台上元信息线索（如项目更新频率、支持者数量等）
通过中心与边缘两条路径显著影响潜在支持者参与决策的证据——
即便这些线索本身并不伴随任何即时经济回报\cite{qiu2025crowdfunding}。
该机制表明，信息本身可作为独立于货币的激励因素发挥作用，
核心路径在于通过改变受众对参与价值的预期评估来引导行为。

知识贡献平台的研究同样提供了重要证据。Toubia 和 Stephen（2013）区分了
“内在效用”（源自贡献内容本身的满足）与”形象效用”
（源自贡献被他人看见和认可的满足），
指出形象效用在很大程度上取决于平台所传递的”贡献受重视”的社会信号\cite{toubia2013intrinsic}。
平台通过公告、排行榜或特殊徽章等方式显式传递此类信号时，
用户参与意愿受到显著影响——尽管这些信号并未带来任何直接金钱收益。
Du（2023）在奖励型众筹场景中的研究进一步证实，
项目信息披露的详细程度直接影响支持者的投资意愿\cite{du2023impact}；
Bai 等人（2024）在金融信息披露领域也注意到，
平台提供的投资者信息透明化机制能显著改善资产定价效率\cite{bai2024platform}。
以上证据共同表明，信息披露能够通过激活受众期望、修正其对奖励可得性的判断，
产生独立且可测量的行为效应，为本研究核心假设提供了实证支撑。

从比较视角看，Brink 的捐赠信息披露与上述研究情境兼有共性与差异。
共性方面，信息的公开传播先于或独立于实际经济转移，
且信息受众（所有关注 Brink 的开发者）远大于实际经济受益者（获得资助的4名开发者）。
差异方面，开源开发者贡献的回报机制更为复杂：
他们面临的不仅是”是否参与”的二元决策，还有”投入多少”的连续决策，
且该决策嵌入在多年社区参与历史中，受声誉、职业发展和内在动机的交织影响。
这意味着 Brink 信息披露的效应可能比众筹场景中的进展公告效应更为细微、更具情境性。
但分析价值恰在于此：若在如此复杂的动机结构中，
信息披露仍能产生可识别的增量效应，则说明信息因素对开源开发者行为具有
超越噪音水平的独立影响力——这一发现将具有重要理论意涵。

\section{期望理论与期望失验理论}

\subsection{期望理论}

期望理论由心理学家 Victor Vroom 于1964年在工作动机研究中
正式提出\cite{vroom1964work}，是组织行为学和管理学领域引用最为广泛的动机理论之一。
该理论的核心命题在于：个体行为动机并非由需求本身决定，
而取决于其对行动结果的认知预期——具体地说，个体愿意付出努力的程度
由三个认知因素的乘积决定。
其一为期望值（Expectancy），即个体相信自身努力
能达到预期绩效水平的主观概率；
其二为工具性（Instrumentality），即个体相信达到特定绩效
能带来相应奖励的主观概率；
其三为效价（Valence），即个体对特定奖励结果的主观价值评估。
任何一个因素趋近于零，整体动机水平便随之崩溃。
这解释了为何在不信任奖励可得性或不重视奖励内容的环境中，
外部激励往往难以奏效。

期望理论最初用以解释雇佣关系中员工的努力水平，
但”努力→绩效→奖励”的因果链逻辑具有高度跨情境适用性，
已被广泛延伸至志愿行为、公民参与和在线贡献等非雇佣情境。
开源社区中，”奖励”的概念扩展至声誉提升、技术满足感和社区认可等多种形式。
而一旦外部货币激励机制进入社区（如 Brink 的资助项目），
货币奖励便成为开发者感知奖励结构的重要组成部分，
期望理论的逻辑也随之获得更为直接的适用空间。

期望理论在在线贡献情境中的适用性已获若干实证研究的初步支持。
众筹领域中，Du（2023）证实项目进展信息更新
能提升潜在支持者对项目最终成功的预期，进而提升其贡献意愿，
这条路径与期望理论的期望值维度高度吻合\cite{du2023impact}。
在线问答社区中，秦敏和梁溯（2017）确认，
社区的用户识别机制（如通过积分、徽章等方式公开标识贡献水平）
能显著提升用户对工具性的感知，进而激励持续的知识贡献行为\cite{qin2017user}。
货币奖励机制（如平台积分或徽章）
被发现主要通过提升工具性感知（高质量回答确实会获得奖励）激励知识贡献，
而非仅凭奖励的直接经济价值发挥作用。
上述研究共同指向一个结论：奖励信息的可见性与可信性——
即潜在贡献者能否清楚地看到并相信”努力→奖励”的实现路径——
对于外部激励能否有效作用于行为，其重要性与激励本身同等甚至更为关键。
Brink 的信息披露实践从信息可见性角度为开发者提供了该路径的具体依据，
其作用机制正是通过修正期望理论三要素的感知水平得以实现。

就本研究的具体情境而言，Brink 的非定向披露向整个开发者社区传递
“资金存在、奖励可得”的信号，同时作用于期望理论的两个维度：
捐款金额的公开提升了开发者对资助项目规模与持续性的感知，
从而增强了奖励的效价；
资金的公开存在又使开发者对”自己也有可能被资助”
形成更积极的预期，提升了期望值。
定向披露则通过揭示具体受资助者及其资助理由，
向未入选开发者提供了”优秀贡献确实会被认可和奖励”的证据，
强化了绩效与奖励之间的因果连接，提升了工具性。
基于以上分析，本研究预期两种披露均正向影响开发者的贡献行为，
但所激活的认知路径有所不同。

\subsection{期望失验理论}

期望失验理论由消费者行为研究者
Richard Oliver 于1980年提出，旨在解释消费者满意度的形成机制\cite{oliver1980cognitive}。
该理论主张：个体满意度并非取决于实际结果的绝对水平，
而取决于实际感知与事先期望之间的比较关系。
实际结果超过期望时，个体产生正向失验，体验到满意乃至喜悦；
实际结果低于期望时，个体产生负向失验，体验到失望乃至抱怨；
结果与期望相符时，期望得以确认，个体保持中性或满足状态。
该框架最初在消费者满意度研究中取得广泛实证支持，
后延伸应用于信息系统使用后评价\cite{abrate2021relationship}等多种情境，
在解释重复使用意愿、持续贡献行为等方面展现了稳健的解释力。

期望失验理论的核心贡献在于阐明了期望水平作为比较基准的关键作用：
同样的实际结果，置于不同期望背景下可能引发截然相反的情绪反应和后续行为。
对未被 Brink 资助的开发者而言，
非定向披露所激活的”我也可能被资助”构成了一个比较基准。
定向披露随后公布受资助者名单时，这批开发者会将自身状况与公告结果加以比较——
绝大多数未入选者由此产生明确的负向失验：
期望未能实现，伴随而来的失望情绪可能削弱其参与意愿，
尤其削弱其对后续非定向披露所激活期望的响应强度。
这种抑制效应，正是本研究预测定向披露对非定向披露具有负向调节效应的核心理论依据。

需指出的是，期望失验效应的强度受初始期望水平的调节。
Abrate 等人（2021）对价格—体验不匹配的研究表明，
初始期望越高、结果的负向偏离越大，失望情绪越强烈；
初始期望本就较低时，负向失验的冲击也相应较弱\cite{abrate2021relationship}。
本研究情境中，不同贡献水平的开发者对获得 Brink 资助的初始期望存在差异：
重度贡献者可能认为自己”更有资格”，初始期望因此更高，
负向失验后的动机削弱也更为明显；
轻度贡献者则可能将期望预设在较低水平，受定向披露影响相对有限。
该理论预期为后续的开发者异质性分析提供了重要指引。

\subsection{两种理论的整合：序贯披露的复合效应}

期望理论与期望失验理论在 Brink 两阶段披露的时间序列中形成自然互补，
共同描述了一条完整的认知动态链条。
非定向披露率先激活期望，在全体开发者中提升效价与期望值，初步促进贡献。
定向披露随后将这一过程推向结果阶段，效应在两个层面同时展开：
就全体开发者而言，定向披露通过展示成功案例增强工具性，产生正向主效应；
就大多数未入选者而言，定向披露又触发负向期望失验，
削弱了其对后续非定向披露所激活期望的响应强度，由此产生负向调节效应。
“整体正向、交互负向”的复合结构
在同一研究情境中为两种理论提供了同步检验的机会，
亦是本研究在单一数据框架内整合两种理论的逻辑基础。

需特别指出的是，期望理论与期望失验理论的整合在信息系统与在线平台研究中
已有先例，但具体应用场景各有侧重。
消费者技术接受领域中，Oliver（1980）的失验框架经 Bhattacherjee（2001）\cite{bhattacherjee2001understanding}
扩展为IS持续使用研究的核心模型——“确认模型（Expectation-Confirmation Model）”，
主张用户在初次使用后会将实际体验与使用前期望进行比对，
确认或失验这一认知过程直接决定其持续使用意愿。
本研究的对象虽为开发者活动而非技术产品的持续使用，
底层逻辑却高度相似：非定向披露激活期望（类比使用前的期望形成），
定向披露公布受益人（类比使用后的实际结果感知），
期望与结果的落差（类比确认/失验）最终影响对后续披露的响应强度。
结构上的这种相似性不仅为本研究的理论借用提供了依据，
也意味着研究发现可与信息系统持续使用研究领域展开跨领域对话，
在理论贡献的辐射范围上具有一定延展性。
另需注意的是，期望水平本身在本研究情境中并不固定，
而是随非定向披露次数的积累和社区对 Brink 运作模式的了解而动态调整。
随着开发者逐渐认识到获得 Brink 资助的实际门槛，
初始期望的高估程度会随时间递减，负向失验效应也可能相应减弱。
这一动态性质提示，研究须在面板模型中对时间趋势加以控制，
以免将随时间演变的期望调整过程错误归因于特定披露事件的效应。

\section{情感与信息处理}

\subsection{情感即信息理论}

情感即信息理论（Affect-as-Information Theory）由 Schwarz 和 Clore（1983）
在一系列以情绪和判断为主题的经典实验中提出\cite{schwarz1983moods}。
该理论的核心主张是：个体作出判断和决策时，
会将当前感受到的情感状态作为关于处境质量的信息来源加以利用，
由此影响判断和行动的方向。
情感状态不仅是对事件的情绪反应，其本身携带着关于环境状态的认知内容——
积极情感传递”当前状况良好、无需采取额外行动”的信号，
促使个体倾向于维持现状或从事探索性活动；
负面情感传递”当前存在问题或威胁、需要警惕与应对”的信号，
促使个体转向系统性分析和问题解决导向的行为。

杨{\notoserifjk 洸}（2020）通过社交媒体的实证研究确认，
网络情感传染遵循特定的线索影响机制：情感信号在群体间的扩散速度
远超中性信息，且负面情感的传染效应显著强于正面情感\cite{yang2020sentiment}。
开源社区情境中，Issue 讨论区是社区成员集体情感最为集中的体现场所。
项目遭遇技术瓶颈、安全漏洞或核心开发者之间的意见分歧时，
Issue 讨论的情感极性会出现明显的负面偏移——
语气趋于焦虑、沮丧或激烈，隐含着”项目面临困难、需要更多投入”的集体信号。
依据情感即信息理论，观察到这种集体负面情感的开发者
会将其解读为”项目处于需要行动的危机状态”，
进而激发增加代码活动（通过提交修复代码直接应对问题）
或增加讨论参与（通过协商寻求解决方案）的动机。
本研究中，情感即信息理论构成了社区负面情感
正向影响代码活动和讨论活动这一预期（H4、H5）的核心理论依据。

需注意的是，情感信号所触发的行为效应并非当期即刻呈现，
而是在情感信号被个体感知、加工后，经一定认知消化期方才影响行为决策。
换言之，情感变化与行为变化之间存在滞后关系。
这一预期与向量自回归模型所捕捉的动态传导机制天然契合，
使情感即信息理论框架与本研究采用的时间序列方法论形成相互支撑。

\subsection{目标手段替代理论}

目标手段替代理论（Goal-Means Substitution Theory）源自目标系统理论
（Goal Systems Theory），由 Kruglanski 等人（2002）系统提出\cite{kruglanski2002theory}。
目标系统理论认为目标与实现目标的手段之间形成关联网络：
同一目标可通过多种手段实现时，若某一手段已被充分激活和使用，
与该目标关联的其他替代手段激活程度会相应降低。
内在机制在于：目标激活（即个体对实现目标的需要感）驱动手段选择；
当某一手段的执行使目标实现程度得到充分满足，
对替代手段的需求感随之减弱。
加之个体时间与精力资源有限，已高度激活的手段会占用可用于替代手段的注意力资源，
进一步压缩其使用空间。

开源社区协作情境中，代码活动（提交代码修复问题）
与讨论活动（通过交流协商解决方案）可被理解为
服务于”解决项目问题”或”维持项目进展”这一共同目标的两种替代性手段。
社区成员已大量提交代码修复、有效推进问题解决时，
“解决项目问题”的目标通过代码手段得到了相当程度的满足，
对讨论这一替代手段的需求感相应降低。
据此，本研究预期代码活动的增加会在滞后一定时期后
对讨论活动产生负向替代效应（H6）。
上述分析还暗示替代效应的强度具有情境依赖性：
代码质量高、修复进展顺利的时期，替代效应更为明显；
代码进展受阻、问题悬而未决的时期，
讨论的必要性不会因代码活动增加而减少，替代效应随之减弱。

\subsection{开源社区中的情感研究综述}

随着自然语言处理技术日趋成熟，情感分析（Sentiment Analysis）
逐渐成为开源社区研究的重要工具。
Guzzi 等人（2013）分析 Apache 和 Mozilla 项目邮件列表后指出，
开发者之间的情感沟通（如表达感谢、沮丧、激动）
在技术决策讨论中占有相当比例，
这些情感内容构成了社区社会资本与协作质量的重要信号\cite{guzzi2013}。
Schroter 等人（2010）通过 Bug 报告质量的分析证实，
信息充分的问题报告（如包含堆栈跟踪信息）更可能获得开发者的有效修复，
揭示了信息质量对技术协作效率的实质性影响\cite{schroter2010}。
Ortu 等人（2015）在 MSR 会议上发表的研究则呈现了一个颇具张力的结果：
Issue 讨论中表达较强负面情感的开发者反而展现出更高的问题修复效率，
暗示负面情感可能通过激活问题导向的行为动机对生产力产生正向影响\cite{ortu2015bullies}。
Novielli 等人（2018）系统评估了多种情感分析工具在软件工程文本上的表现，
确认通用情感词典在技术讨论场景中的准确性显著低于预期，
并呼吁研究者应用情感分析工具时须充分考虑领域适配性\cite{novielli2018sentiment}。
Xie 和 Zhang（2021）在 CSCWD 会议上报告的研究直接考察了贡献者情感
对知识贡献意愿的影响，注意到负面情感体验虽降低了个体的即时满意度，
却在短期内激发了更强的问题解决动机\cite{xie2021influence}。
Li 等人（2021）开发了一套负面情感事件监测框架，
可从开源项目的 Issue 讨论中自动检测情感异常事件，
为社区管理者的干预决策提供数据驱动的支持\cite{li2021monitoring}。
上述研究初步确立了情感在开源社区中所扮演的信息传导角色，
但普遍采用截面分析或短期事件研究，难以捕捉情感变化与行为变化之间的动态时间关系。

更深层的局限在于，现有研究对情感与不同类型开发活动
（代码活动与讨论活动）之间差异化影响路径关注不足，
更缺乏对代码活动与讨论活动之间动态替代关系的系统考察。
具体而言：情感如何分别影响代码活动和讨论参与？
代码活动的增加是否会随后压缩讨论活动的空间？
这一涉及情感与行为动态传导完整链条的问题，
在开源社区研究中尚属空白。

现有开源情感研究在方法论层面同样存在不容忽视的局限。
早期研究大多使用基于词典的情感分析工具（如 SentiStrength、LIWC），
此类工具在通用文本上表现尚可，
但高度技术化的开源讨论文本中，大量领域专属表达（如代码片段、版本号、
错误堆栈信息）会对情感评分产生系统性干扰，准确性因此受限。
近年来，基于预训练语言模型的情感分析方法（如 BERT\footnote{BERT（Bidirectional Encoder Representations from Transformers）由 Google 于2018年提出，是一种基于 Transformer 架构的双向预训练语言模型，在多项自然语言处理任务上取得了突破性进展。} 及其衍生模型）
在技术文本理解上取得了显著进步。
Novielli 等人（2021）对多种现成的软件工程专用情感分析工具展开跨平台评估，
报告工具性能在不同数据源之间差异显著，
并呼吁研究者选择工具时考虑目标平台的特性\cite{novielli2021assessment}。
Cassee 等人（2024）进一步验证了基于 Transformer 架构的情感分析模型
在软件工程文本上的表现，
确认元标记化（Meta-Tokenization）策略能有效提升模型对技术术语的处理能力\cite{cassee2024transformers}。
李合龙等（2023）在金融市场文本情绪研究综述中系统梳理了情绪度量方法的演变路径——
从基于词典的方法到深度学习模型——为跨领域情感分析方法的选择
提供了有价值的参考框架\cite{li2023sentiment}。
但这些方法在开源社区情感分析中的系统性评估仍相对有限。
如何选择适合技术文本的情感分析工具、如何处理专业领域词汇对通用情感模型的干扰，
本身即为开源社区情感研究面临的方法论挑战之一。
本研究将在数据与方法章节对此作出专门说明，
并通过替换情感工具的稳健性检验验证主要结论对工具选择的不敏感性。
借助包含情感、代码活动和讨论活动的 VAR 模型，
本研究在方法论和研究问题两个层面均对既有研究作出实质性推进。

\section{开源社区中的活动类型与相互关系}

GitHub 作为全球最大的开源代码托管平台，
提供了高度结构化的开发活动记录。
这些记录构成了量化研究开源社区协作行为的核心数据资产\cite{alami2024sustainability}。
GitHub 平台上开发者可参与多种形式的活动，
不同活动类型在与项目最终产出的直接关联程度、参与门槛以及
对外部信号的响应机制上均存在实质性差异，有必要分别对待。

代码活动（Code Activity）涵盖开发者对项目代码库的直接操作行为，
包括 PushEvent（PE）、CreateEvent（CE）、DeleteEvent（DE）、CommitCommentEvent（CCE）和 PullRequestEvent（PRE）五种事件类型。
代码活动要求开发者具备足够的技术储备和对项目状态的深度了解，
通常需要投入相对连续且高度集中的时间，是反映开发者核心产出的最直接指标。
审查活动（Review Activity）是开发者对他人提交代码进行质量评估的协作行为，
包括 PullRequestReviewEvent（PRRE）和 PullRequestReviewCommentEvent（PRRCE）两种事件类型。
代码审查要求审查者深入理解项目规范和技术标准，
是维护代码质量、防止技术债务累积的关键机制，
在大型开源项目中也往往是新贡献能否被合并的决定性环节。
讨论活动（Discussion Activity）以 IssuesEvent（ISE）度量，
涵盖问题报告、需求提案、功能讨论和技术争辩等多种形式，
参与门槛相对较低，任何具备基本技术背景的社区成员均可参与，
是社区协商和集体决策的主要渠道。
需要说明的是，IssueCommentEvent（ICE）未被纳入讨论活动的度量，
原因在于 ICE 既包含对 Issue 的评论，也包含对 Pull Request 的评论，
其行为含义与审查活动存在重叠，纳入将导致活动分类边界模糊，故予以剔除。

三类活动不仅功能有别，
在对外部信息信号的响应机制上也可能呈现差异。
代码活动所需时间成本最高，其变化往往反映深层动机改变，
响应外部信号的时间窗口相对较长；
讨论活动时间成本较低，对外部刺激的响应更快、弹性更大；
审查活动则介于两者之间。
对三类活动分别建模，有助于揭示不同激励机制在不同活动层面的差异化影响，
从而对开源开发者行为结构作出更为精细的刻画\cite{schroter2010}。
顾美玲等（2019）对开放式创新社区的研究亦指出，
社区治理机制对不同类型知识贡献行为具有差异化影响，
过度简化的合并指标会掩盖重要的异质性信息\cite{gu2019governance}。

另一方面，三类活动并非相互独立，而是通过开发者有限的时间与精力资源相互关联。
代码活动与讨论活动均服务于”推进项目解决问题”的共同目标，
但所需资源和关注焦点不同——代码活动要求高度专注的技术投入，
讨论活动则更依赖对社区动态的关注和沟通意愿。
本研究引入目标手段替代理论，对代码活动与讨论活动之间可能存在的替代关系
提出明确的理论预期，并通过 VAR 模型对这一动态路径进行实证检验，
从而将活动类型的理论区分与动态因果分析有机结合起来。

从研究设计角度看，三类活动的分开测量还具有规避”合并谬误”的方法论意义。
若将所有开发活动合并为单一指标，
激励信号对不同活动类型的差异化效应可能被掩盖。
举例而言，非定向披露可能更强烈地激励门槛较低的讨论参与，
对需大量时间投入的代码活动影响有限；
定向披露所展示的榜样效应则可能更强烈地激励能直接展示技术能力的代码活动。
类似地，社区情感对讨论活动的即时影响可能强于对代码活动的影响——
讨论对情感信号的响应更为迅速，代码提交往往需要较长准备周期。
对三类活动分别建模，使本研究得以更精细地捕捉激励机制的作用边界，
避免因指标聚合遮蔽有价值的异质性信息。
这种细粒度的活动分类也赋予研究发现更具体的实践参考价值：
管理者可据此判断哪种类型的激励策略
对激活哪种类型的活动最为有效。

\section{本章小结}

本章围绕两项研究的理论需求，系统梳理了五个相互关联的理论与文献板块。
开源开发者激励文献揭示了内在动机主导下货币激励可能产生挤出效应的复杂格局，
同时指出现有研究忽视了信息披露本身的激励价值。
自愿信息披露理论和信号理论为理解 Brink 两阶段披露机制的信号性质提供了概念工具；
期望理论和期望失验理论则为预测两种披露的主效应及其交互效应
提供了心理机制层面的解释框架。
情感即信息理论和目标手段替代理论为社区情感影响开发活动、
进而触发代码—讨论替代关系提供了理论依据；
对开源社区活动类型及其相互关系的梳理，
确立了本研究区分三类开发者活动分别建模的方法论合理性。
上述理论框架共同构成后续两项实证研究的逻辑基础，
并在各章假设推导与结果解读中得到具体应用。


